\section{Bidirectional, Agentic, Multimodal, and Temporal Systems}
\label{sec:bidirectional}

The families above treat construction and grounded generation as separate directions. The frontier is increasingly bidirectional. A model extracts or proposes graph content, the graph guides later retrieval, and new evidence can trigger another construction cycle. The loop is attractive because biomedical knowledge is incomplete and moves fast. It is also where a useful retrieval component quietly becomes a knowledge-management system. Errors now persist across sessions. Generated claims can pass for external evidence. Updates change behavior without a release. The real question is not whether the model and graph exchange information. It is whether the exchange keeps track of what came from where.

\subsection{Bidirectional co-augmentation}

DALK is an early biomedical example. It uses a model to build an evolving Alzheimer's-disease graph from the literature, then brings selected graph evidence back into inference through coarse-to-fine sampling~\cite{li2024dalk}. This addresses a genuine limitation of static graphs, since specialized evidence is often too recent or too granular to be there. It also exposes an ambiguity in the word ``bidirectional,'' which covers two very different claims. One is an engineering claim: construction output becomes retrieval input. The other is an epistemic claim: the resulting graph is reliable enough to constrain answers. Only the first follows automatically. The second needs source-span grounding, predicate validation, versioned snapshots, and testing against evidence the same model did not produce.

Generate-on-Graph, though evaluated outside biomedicine, offers a useful abstraction. Its loop lets the model explore the graph and synthesize missing triples when the graph alone cannot complete a path~\cite{xu2024gog}. Synthesis like this can raise recall. But a generated triple is an inference, not retrieved evidence. In a clinical deployment it belongs in a scratch graph scoped to the current request. It may support a hypothesis or shape a search query. It must not inherit the status of a curated edge merely because it completed a plausible path.

KGARevion runs a safer asymmetric loop. The model proposes candidate triples, and the graph verifies and filters them before an answer is written~\cite{su2025kgarevion}. The model supplies recall, the graph keeps precision. This works well when the graph is authoritative in its domain. One thing stays ambiguous, though: ``not found'' can mean the claim is false, or poorly normalized, or out of scope, or simply absent from this release. Verification should say which. A typed verdict, supported or contradicted or unresolved or out of scope, is more useful than collapsing absence into falsity.

Table~\ref{tab:loops} sorts these systems by what actually crosses the feedback boundary. The distinction heads off a common error: a system can be bidirectional at inference without being allowed to write to the institutional graph.

\begin{table}[!t]
\centering
\caption{\textbf{Bidirectional systems, distinguished by what actually crosses back into the graph.} ``Bidirectional'' is often used as a single label, but the loops differ in the one respect that matters for safety: how long the model's output persists and what authority it acquires. A hypothetical edge invented to complete a reasoning path should live only for the current request. A candidate relation extracted from a paper may sit in a staging area awaiting review. Only a curator-approved assertion belongs in the graph that later queries treat as authoritative. The rightmost column is therefore a design requirement rather than a description: a system that lets a reasoning-loop artifact acquire maintenance-loop persistence is the circular self-validation failure of Section~\ref{sec:failure}.}
\label{tab:loops}
\footnotesize
\begin{tabularx}{\columnwidth}{@{}L{1.7cm} Y L{1.9cm}@{}}
\toprule
\textbf{Loop type} & \textbf{LLM-to-graph / graph-to-LLM artifact} & \textbf{Persistence} \\
\midrule
Reasoning & Hypothetical edge; neighbor/path evidence & Request-scoped \\
Retrieval & Query reformulation; source-linked subgraph & Logged, not asserted \\
Curation & Candidate relation; constraint and curator feedback & Candidate store \\
Learning & Synthetic instruction/reward; graph supervision & Versioned artifact \\
Maintenance & Proposed addition/correction; existing claims and provenance & Quarantined until approval \\
\bottomrule
\end{tabularx}
\end{table}

\subsection{When is a system genuinely agentic?}

The word \emph{agentic} gets attached to any multistage prompt chain. We reserve it for something narrower: a controller that watches an evolving state, picks among available actions, judges what comes back, and decides whether to continue, revise, abstain, or stop. A fixed sequence of link, retrieve, generate is a pipeline, even when every stage calls a model. An agent might broaden a subgraph, fetch the source article, issue a structured query, test a contradiction, or ask for human review, depending on what it found. The distinction is not pedantry. A pipeline's failure modes can be listed in advance. An agent's depend on trajectories that only exist at run time.

\subsection{The tool surface over a medical graph}

In practice an agent is defined by the tools it can call, which makes the tool surface the main design artifact. Six operations are worth separating, because each carries a different risk.

\emph{Concept resolution} maps a mention to identifiers. It should return ranked candidates with confidence, not one answer, because an early commitment to the wrong identifier constrains everything downstream. \emph{Neighborhood expansion} retrieves typed neighbors under an explicit hop and degree budget; without a budget, one hub concept exhausts the context window. \emph{Path query} finds routes between anchor concepts under relation-type constraints. \emph{Structured query execution} runs a generated SPARQL or Cypher statement, read-only, against an allowlisted grammar. \emph{Source retrieval} fetches the documents behind an edge, which is how an agent recovers the dose, population, and study design a triple cannot carry. \emph{Verification} checks a claim against the graph.

Two rules follow from how these systems fail. Verification should return a typed verdict rather than a boolean, because collapsing absence into falsity turns a graph-coverage limit into an apparent factual finding~\cite{su2025kgarevion}. And write access belongs on its own tool with its own authorization. An agent that can propose an edge and an agent that can assert one are different risk classes, which Section~\ref{sec:failure} takes up.

\subsection{Planning, control, and stopping}

General graph agents supply traversal patterns that biomedical systems inherit. Think-on-Graph runs an iterative beam search, letting the model pick promising relations and entities~\cite{sun2024tog}. Reasoning on Graphs separates planning from execution, generating relation-path plans grounded in the graph before any answer~\cite{luo2024rog}. That separation is the more transferable idea. A plan written as a relation path can be checked against the schema before retrieval runs, which catches a whole class of error cheaply. Think-on-Graph 2.0 alternates graph and document retrieval, so structure deepens the document search while text refines the graph exploration~\cite{ma2025tog2}. The alternation suits biomedicine well. The graph finds the relational route; the documents supply the dose, population, and study design that no triple carries.

Control is where biomedical agents usually go wrong. Every extra action widens the attack surface and leaves more room for confirmation bias. An agent that keeps expanding until it finds support will always find some. Three controls belong in any implementation. Cap the tool calls and the cumulative retrieved context. Require a search for disconfirming evidence, not just corroboration. Expose the stopping criterion, whether that is evidence sufficiency, a contradiction threshold, a budget, or an abstention. HyKGE shows the hazard from the retrieval side, since a provisional hypothesis used to guide exploration can steer the search toward paths that agree with it~\cite{jiang2025hykge}. An agent that writes its own hypotheses has the same problem and more iterations in which to compound it.

\subsection{Multi-agent decomposition}

Splitting a task across specialized agents has the strongest clinical evidence behind it. The oncology trial-matching platform separated extraction from eligibility reasoning. Collapsing the two into one call lowered extraction fidelity and raised unresolved conflicts, dropping F1 from 0.82 to 0.78~\cite{loaizabonilla2026trialmatch}. That is a rare direct measurement of a decomposition benefit. The mechanism is worth naming: keeping extraction and reasoning separate keeps each stage inspectable, and stops an error in one from being laundered through the other.

The same study shows the complementary control. Its final eligibility verdict came from deterministic rules, not a generator, with clinician review on top. CrossDDI makes the same argument for drug interactions, taking the model out of arbitration entirely and issuing a positive verdict only on explicit supporting evidence, which cut variation across three backbones~\cite{wang2026crossddi}. For any decision expressible as a rule, the defensible design is agents for gathering and a deterministic component for deciding. It costs some coverage. It buys traceability and stability across model versions.

\subsection{Agents that maintain the graph}

AMG-RAG is the clearest case of an agent whose job is the graph rather than the answer. It automates construction and continuous updating, and retrieves current graph evidence for question answering~\cite{rezaei2025amgrag}. MedKGent scales the ingestion side, using two agents to extract and integrate relations from over ten million PubMed abstracts while resolving conflicting claims~\cite{zhang2025medkgent}.

Benchmark accuracy is the wrong endpoint for both. A wrong answer is a visible, bounded error. A wrong persistent edge is neither. It contaminates every later retrieval, training set, explanation, and decision that touches it, silently. Update evaluation needs its own metrics: candidate-edge precision by relation type, conflict-resolution accuracy when sources disagree, time from published correction to graph correction, how far a bad edge propagates, and whether rollback actually works. None of these is reported anywhere in the current literature. That gap matters more than any remaining headroom on question-answering accuracy.

\subsection{What an auditable agent trace contains}

Reproducing an agent run takes more than the prompt. An auditable trace should record the initial request and policy context, each state representation, the chosen action alongside the alternatives that were available, exact tool inputs and outputs, identifiers and versions for the model, graph, ontology, retriever, and prompts, source timestamps, evidence retrieved and rejected, confidence or abstention decisions, and the terminal reason. This is not equivalent to publishing hidden chain-of-thought, which may be incomplete, unstable across runs, or disclosive. A concise event log captures the externally meaningful decisions without storing unrestricted internal reasoning, and it is what makes a failure attributable after the fact. Reproducibility additionally needs a snapshot mode, so the same request can be replayed against frozen dependencies even while the production graph continues to update.

\subsection{Multimodal graphs as an alignment layer}

Biomedical entities are described by text, images, molecular structures, sequences, physiological signals, and clinical events, and a graph can align them. The phrase covers three distinct designs, which should not be evaluated interchangeably: encoders aligned through graph relations, as in BioBridge~\cite{wang2024biobridge} and the sequence-augmented PrimeKG++ resource~\cite{dang2025biomedkg}; multimodal evidence stored as graph attributes and retrieved at inference; and graph-derived supervision used to train a multimodal model, as in MedKInstruct~\cite{wu2026medkinstruct}. The first two are graph--foundation-model integrations rather than generative systems, so their endpoint is cross-modal retrieval quality and inductive performance on unseen entities, not answer faithfulness. The third relocates the graph from context to supervision, which improves compositional reasoning and makes leakage harder to see: a test answer or near-duplicate caption present during instruction synthesis lets the model reproduce graph-to-label mappings without attending to the image. Image-counterfactual tests and entity-disjoint splits are the minimum defense.

Physiological signals are the newest frontier. EEG-MedRAG organizes domain knowledge, patient cases, and a repository into a three-layer hypergraph for joint semantic and temporal retrieval~\cite{wang2025eegmedrag}, which is a reasonable answer to the n-ary clinical relations that plain triples flatten, and which will need prospective testing across devices and institutions before it means much clinically. Across all three designs, a multimodal system should report how it behaves when a modality is missing, how calibrated it stays when modalities disagree, and whether a given explanation rests on the patient signal or on generic graph knowledge.

\subsection{Temporal knowledge is more than a fresh graph}

Biomedical claims run on four clocks, and conflating them is a common design error. Transaction time records when an edge entered a database, publication time when the supporting evidence appeared, valid time when a recommendation actually applies, and patient-event time when something happened to a specific person. A graph rebuilt every night has excellent transaction-time resolution and can still be temporally naive, because none of that tells you whether a recommendation was superseded.

MedKGent shows the ingestion side is tractable, using two agents to extract and integrate relations from more than ten million PubMed abstracts spanning several decades and building a daily series of graphs while resolving conflicting claims~\cite{zhang2025medkgent}. A series of snapshots, though, captures publication and transaction history rather than clinical time. What clinical reasoning needs are statements of the form ``this symptom typically precedes that event,'' or ``this therapy was contraindicated during that interval.'' ChronoMedKG is the first resource to attempt them at scale, reporting 460{,}497 evidence-linked triples over 13{,}431 diseases with onset windows and traceable support, alongside a temporal question benchmark with static controls; performance drops from static to temporal questions and graph retrieval recovers only part of it~\cite{ahmed2026chronomedkg}. The work was a preprint at the time of writing, so its numbers are indicative rather than settled, but the direction is hard to dismiss.

The evaluation requirement that follows is an as-of boundary. For a question asked at time $t$, every edge, document, index, adaptation set, and cached summary must be restricted to what existed by $t$. Random edge splits leak future knowledge even when the visible prompt contains only older facts, and a contemporary model leaks more. Strong evaluations add evolving-guideline cases, retractions, and adversarial conflicts between older high-volume evidence and newer authoritative guidance. A system should also be explicit about which question it answers: what is believed now, what was known then, or what applied to this patient at the time. These are not interchangeable.

\subsection{Preventing circular self-validation}

Closed-loop systems create a distinctive failure mode. Suppose model $M$ proposes claim $c$, a graph updater stores $c$, and a later instance of $M$ retrieves that edge as ``external'' support. The loop has converted one model output into apparently independent evidence. Self-consistency, repeated extraction, or agreement among closely related models can amplify the effect without adding a genuinely independent source. The resulting citation can be syntactically correct and epistemically circular.

Four controls are mandatory. Maintain separate asserted, candidate, and ephemeral graph zones, and let only asserted edges serve as authoritative grounding. Require every persistent claim to point to a resolvable primary or curated source span, with model identity recorded as an extractor rather than an evidence source. Require independent validation before promotion, by a curator, a deterministic database check, an independently sourced document, or a suitably independent model with calibrated disagreement handling. Make promotion, supersession, retraction, and rollback explicit, permissioned events. Agent tools should hold least privilege: retrieval agents get read access, extraction agents may write candidates, and only an authorized release process may write asserted edges.

The same separation must hold through generation. A model-generated summary is a derived artifact, not a primary source; a graph-consistent statement is not necessarily entailed by the cited document; and a source-supported sentence can still be clinically inapplicable because of a population or temporal mismatch. As the sentence-level analysis in VERICITE shows, supplying relevant documents does not by itself ensure faithful claim-citation alignment~\cite{ma2026vericite}. Trustworthy closed-loop systems must preserve lineage from answer claim to graph assertion to source passage, expose uncertainty and conflicts, and abstain when the chain is incomplete. Feedback may improve coverage, but only independently governed evidence may increase authority.
